% ==========================================
% SECCIÓN 1: MARCO TEÓRICO
% ==========================================
% TIP: Incluye aquí los conceptos, definiciones y antecedentes
% necesarios para entender la solución de la práctica.
% ==========================================

\chapter{Marco teórico}

% INSTRUCCIONES GENERALES:
% - Incluir un parrafo introductorio del capitulo antes de la primera seccion.
% - Redactar en prosa (evitar formato de glosario y vietas).
% - Incorporar citas distribuidas en todo el capitulo (IEEE con \cite{clave}).
% - No insertar figuras/tablas antes de texto explicativo.
% - Meta de estilo: 90% prosa y 10% apoyo visual/tablas.

% RÚBRICA DE EVALUACIÓN:
% # Marco teórico

% Aquí se coloca la investigación de los conceptos necesarios y suficientes que se requieren para comprender el desarrollo de la práctica o proyecto.

% Atender las siguientes indicaciones que se tomarán en cuenta en la rúbrica de evaluación:

% * No habrá tolerancia a las faltas de ortografía
% * No habrá tolerancia a la práctica del plagio (copy-paste)
% * Se deben utilizar citas de forma textual o usando paráfrasis
% * El 90% de lo contenido en esta sección debe estar escrito en prosa, dejando el 10% para tablas, numeraciones, incisos y viñetas
% * Esta sección puede ocupar las páginas que se necesiten

% Si se van a ocupar imágenes se deben mencionar en el texto, por ejemplo, en la figura 1 podemos ver la ejecución de un programa en Lenguaje C.
        

% CHECKLIST RAPIDO:
% - Cada seccion debe tener al menos 2 parrafos completos.
% - Evitar copiar texto literal sin cita.
% - Cerrar el capitulo con un parrafo de transicion hacia desarrollo.






